% Chapter 1 - Introduction.
% Main source: docs/vision.md.
% Target length: 4-6 pages.

\chapter{Introduction}
\label{ch:introduccion}

% TODO: opening paragraph. About one page setting the scene: what
% UnivOrch is, why it is built, what problem it solves. It must be
% readable by a tribunal unfamiliar with the project and still convey
% the motivation.

\section{Motivation}
\label{sec:motivacion}

% TODO: draft based on the following ideas:
% - The EI UMA teaching lab combines physical and virtual machines.
% - Each course builds its own pyramid of VMs by hand: tedious and
%   error-prone.
% - Existing orchestrators (vSphere, vRA, OpenStack...) do not cover
%   the teaching case of "one course, N students, one template, IPs
%   from a pool, snapshots, recycling".
% - The supervisor had been working with esxobjects and yamlinfr,
%   which cover part of the problem but are bound to VMware.
% - Build on that experience and generalise: a hypervisor-agnostic
%   engine plus a teaching-specific layer.

\section{Objectives}
\label{sec:objetivos}

% TODO: write as a numbered list of objectives.
% O1. Design and implement the orchestrator core as a hypervisor-
%     agnostic layer.
% O2. Define a declarative model of the resources to deploy, with
%     inheritance and template reuse.
% O3. Establish a connector contract that lets new hypervisors
%     (VMware, Proxmox) be added without touching the core.
% O4. Expose a REST API as the public boundary of the service.
% O5. Provide a command-line client and package the whole system as
%     a single-line installable Docker container.
% O6. Lay the groundwork (extension points and model) for the pieces
%     that the TFG will deliver next: web interface, real connectors,
%     teaching-specific application.

\section{Report structure}
\label{sec:estructura}

% TODO: write a short paragraph describing each of the 5 chapters and
% the 4 appendices. Pattern:
% "Chapter \ref{ch:estado-arte} reviews the technical context...".
